\section{Terry's comment on the forum}

% Faithful paraphrase of Terence Tao's forum comment,
% posted at 22:50 on 1 November 2025.

\begin{proposition}
The assertion
\[
    f(N)=o(N)
\]
does not hold for all positive integers $N$.
\end{proposition}

\begin{proof}
Write the positive divisors of $n$ in increasing order as
\[
    d_1(n)<\cdots<d_{\tau(n)}(n),
\]
and, for $k\geq 0$, define
\[
    \sigma_k(n)
    :=
    \begin{cases}
    d_1(n)+\cdots+d_{\tau(n)-k}(n),
        & \tau(n)>k,\\
    0,  & \tau(n)\leq k.
    \end{cases}
\]
Thus, $\sigma_k(n)$ is the sum of all divisors of $n$ except
for its $k$ largest divisors. In particular,
$\sigma_0(n)=\sigma(n)$.

We first recall that
\[
\begin{aligned}
    \sum_{n\leq x}\sigma(n)
    &=\sum_{\substack{d,m\geq1\\dm\leq x}}d
      =\sum_{d\leq x}d
        \sum_{m\leq x/d}1  \\
    &\ll \sum_{d\leq x}x
     \ll x^2.
\end{aligned}
\]
We shall prove the stronger estimate
\begin{equation}\label{eq:tao-sigma-k}
    \sum_{n\leq x}\sigma_k(n)
    \ll \frac{x^2}{k^{11/10}},
    \qquad k\geq1.
\end{equation}

Assuming \eqref{eq:tao-sigma-k} for the moment, Markov's
inequality gives
\[
\begin{aligned}
&\#\bigl\{(n,k):n\leq x,\ k\geq1,\
                 \sigma_k(n)\geq C_0x\bigr\}  \\
&\qquad\leq
 \frac{1}{C_0x}
 \sum_{k\geq1}\sum_{n\leq x}\sigma_k(n)
 \ll
 \frac{x}{C_0}\sum_{k\geq1}\frac{1}{k^{11/10}}
 \ll \frac{x}{C_0}.
\end{aligned}
\]
The same estimate for $k=0$ follows from
$\sum_{n\leq x}\sigma(n)\ll x^2$. Consequently,
\begin{equation}\label{eq:tao-count}
    \#\bigl\{N\in[C_0x,2C_0x]:f(N)\leq x\bigr\}
    \ll \frac{x}{C_0}.
\end{equation}
Indeed, if $f(N)\leq x$, then there are $n\leq x$ and
$k\geq0$ such that $N=\sigma_k(n)$.

For fixed sufficiently large $C_0$, estimate
\eqref{eq:tao-count} implies that, for arbitrarily large $x$,
there exist $N\in[C_0x,2C_0x]$ such that $f(N)>x$. For these
integers,
\[
    \frac{f(N)}{N}>\frac{1}{2C_0}.
\]
Hence $f(N)/N$ does not tend to zero.

It remains to establish \eqref{eq:tao-sigma-k}. Consider a
divisor $d$ of $n=dm$ that contributes to $\sigma_k(n)$.
Since $d$ is not among the $k$ largest divisors of $n$,
there are at least $k$ divisors of $dm$ that are at most
$m$, by complementary pairing of divisors.

Put
\[
    R_m(d):=\#\{r\leq m:r\mid d\}.
\]
Every divisor of $dm$ that is at most $m$ can be written as
a product of a divisor of $m$ and a divisor $r\mid d$ with
$r\leq m$. Therefore,
\[
    k\leq \tau(m)R_m(d).
\]
It follows that at least one of the inequalities
\[
    \tau(m)\geq k^{1/10},
    \qquad
    R_m(d)\geq k^{9/10}
\]
must hold.

The contribution from the first possibility is bounded by
\[
\begin{aligned}
    \sum_{m\leq x}
    \left(\frac{\tau(m)}{k^{1/10}}\right)^{100}
    \sum_{d\leq x/m}d
    &\ll
    \frac{x^2}{k^{10}}
    \sum_{m\geq1}\frac{\tau(m)^{100}}{m^2}  \\
    &\ll \frac{x^2}{k^{10}}.
\end{aligned}
\]
Here the final series converges, for instance by the usual
divisor bound $\tau(m)=m^{o(1)}$.

For the second possibility, squaring the divisor count gives
\[
\begin{aligned}
&\sum_{m\leq x}\sum_{d\leq x/m}
  d\left(\frac{R_m(d)}{k^{9/10}}\right)^2  \\
&\qquad=
 \frac{1}{k^{9/5}}
 \sum_{m\leq x}
 \sum_{r_1,r_2\leq m}
 \sum_{\substack{d\leq x/m\\ [r_1,r_2]\mid d}}d  \\
&\qquad\ll
 \frac{x^2}{k^{9/5}}
 \sum_{m\geq1}\sum_{r_1,r_2\leq m}
 \frac{1}{m^2[r_1,r_2]}.
\end{aligned}
\]
Since $[r_1,r_2]\leq m^2$, we have
\[
    \frac{1}{m^2[r_1,r_2]}
    \leq
    \frac{1}{m^{9/5}[r_1,r_2]^{11/10}}.
\]
Moreover,
\[
    \sum_{r_1,r_2\geq1}
    \frac{1}{[r_1,r_2]^{11/10}}
    <\infty.
\]
Indeed, writing $r_1=ga$, $r_2=gb$, with $(a,b)=1$, yields
\[
\begin{aligned}
    \sum_{r_1,r_2\geq1}
    \frac{1}{[r_1,r_2]^{11/10}}
    &=
    \sum_{g\geq1}\frac{1}{g^{11/10}}
    \sum_{\substack{a,b\geq1\$a,b)=1}}
    \frac{1}{(ab)^{11/10}}  \\
    &\leq \zeta(11/10)^3<\infty.
\end{aligned}
\]
Thus, the contribution from the second possibility is
\[
    \ll \frac{x^2}{k^{9/5}}.
\]
Combining the two estimates proves
\[
    \sum_{n\leq x}\sigma_k(n)
    \ll
    \frac{x^2}{k^{10}}+\frac{x^2}{k^{9/5}}
    \ll \frac{x^2}{k^{11/10}},
\]
which completes the proof.
\end{proof}

Heuristics concerning the typical divisor structure of
integers suggest that the left-hand side of
\eqref{eq:tao-sigma-k} might even decay exponentially in $k$.